% 3. Antecedentes
\begin{frame}{De dónde venimos}
  \begin{itemize}
    \item SimAS 1.0 sentó las bases: editor de gramáticas y simulación de analizadores descendentes y ascendentes.
    \item SimAS 2.0 amplió funciones: validación automática, más control sobre tablas y errores, e informes; pero aparecieron problemas de consistencia y portabilidad.
  \end{itemize}
\end{frame}

\begin{frame}{Lecciones aprendidas}
  \begin{itemize}
    \item La docencia necesita una herramienta estable: sin pérdidas de estado ni resultados inconsistentes entre pasos.
    \item La visualización es clave: derivación y árbol ayudan a fijar conceptos, pero deben generarse de forma fiable.
    \item La experiencia debe ser fluida: trabajar con varias vistas/pestañas y generar informes de forma sencilla.
  \end{itemize}
\end{frame}

\begin{frame}{Qué aporta SimAS 3.0 respecto a 2.0}
  \begin{itemize}
    \item Correcciones en transformaciones, conjuntos Primero/Siguiente y construcción de la tabla predictiva LL(1).
    \item Derivación y árbol en tiempo real, incluyendo nodos \(\epsilon\), y gestión clara de errores.
    \item Interfaz con pestañas inteligentes, i18n y ejecución consistente en Windows, macOS y Linux.
  \end{itemize}
\end{frame}


